%%%%%%%%%%%%%%%%%%%%%%%%%%%%%%%%%%%%%%%%%%%%%%%%%%%%%%%%%%%%%%%%%%%%%%%%%%%
%%  chapter6.tex  –  Глава 6: Применение единого защитного фреймворка к
%%                            нейросетевой системе навигации беспилотных
%%                            летательных аппаратов
%%  v11 – Новая глава применения, параллельная Главе 5 (RobustIDPS.ai
%%        как промышленная платформа для NIDS). Демонстрирует предметно-
%%        независимость методов M1-M7 через перенос на UAV-домен.
%%        Только заголовки \section добавлены в оглавление.
%%%%%%%%%%%%%%%%%%%%%%%%%%%%%%%%%%%%%%%%%%%%%%%%%%%%%%%%%%%%%%%%%%%%%%%%%%%
\chapter{Применение программной системы RobustIDPS.ai для повышения робастности моделей компьютерного зрения беспилотных летательных аппаратов}
\label{ch:uav}
В настоящей главе систематически реализуется и валидируется
\emph{система} гибридная СОПВ, формализованная в
Главе~\ref{ch:methods}, применительно к сектору защиты беспилотных
летательных аппаратов (БПЛА), охарактеризованному в
\S\ref{subsubsec:sector_uav}. Глава следует пятизвенной схеме,
заданной указанием научного руководителя:
\emph{(I)~постановка задачи защиты системы БПЛА} формализует
входы, выходы и поверхность атак в условиях радиоэлектронной
борьбы (РЭБ) и состязательного машинного обучения
(\S\ref{sec:uav_overview});
\emph{(II)~адаптация гибридной СОПВ к сектору БПЛА}
систематически описывает трёхуровневую архитектуру
\emph{борт~--- дронопорт~--- облако} (\S\ref{sec:uav_arch}) и
формализует операционный граф БПЛА как экземпляр семёрки
$\mathcal{H}=(\mathcal{V}_t, \mathcal{E}_t, X_t, A_t, f_\theta,
g_\psi, \pi)$ из \S\ref{sec:hybrid_idps_math_model}
(\S\ref{sec:uav_graph});
\emph{(III)~применение разработанных моделей М1--М7 к подсистемам
БПЛА} систематически представляет отображение методов на
конкретные подсистемы (зрительное восприятие, ГНСС-навигация,
автопилотная политика управления) и сертификаты
Липшица--Грёнуолла и рандомизированного сглаживания, применённые
к UAV-задачам
(\S\ref{sec:uav_pipeline}--\S\ref{sec:uav_certs});
\emph{(IV)~программная реализация системы для сектора БПЛА}
описывает референсную реализацию Phase~A на основе библиотеки
\textsf{robustnn-core} с её интеграцией в программную платформу
\textsf{RobustIDPS.ai}
(\S\ref{sec:uav_phase_a}--\S\ref{sec:uav_robustidps_integration});
\emph{(V)~апробация системы} систематически документирует
результаты валидации Phase~A с операционной метрикой
\emph{Mission-Completion-Rate vs.~$J/S$} (характерный целевой
рабочий режим $\textsf{MCR}\ge 0{,}80$ при отношении
помеха/сигнал $J/S \le 20\,$дБ), сравнение с существующими
промышленными решениями и соответствие нормативной базе
(\S\ref{sec:uav_phase_a_results}--\S\ref{sec:uav_regulation}).
Литературный обзор сектора БПЛА и сравнительный контекст
вынесены в~\S\ref{sec:ch1_uav_related} Главы~\ref{ch:literature};
математические основы~--- в Главу~\ref{ch:methods};
предварительная валидация гибридной СОПВ на UAV-задачах в
качестве преамбулы~--- в Главе~\ref{ch:adaptation}.
%%=====================================================================
\section{Угрозы безопасности моделей компьютерного зрения беспилотных летательных аппаратов}
\label{sec:uav_overview}
%%====
Современный БПЛА является нейросетево-зависимой системой на всех
уровнях программного стека: бортовые свёрточные и трансформерные
детекторы обеспечивают зрительное восприятие, рекуррентные и
гауссово-процессные модели~--- оценку ГНСС-целостности, политики
глубокого обучения с подкреплением (PPO, DDPG-PER)~--- внешний контур
автопилота, графовые модели~--- координацию роя и оркестрацию
дронопортов. Каждый из этих компонентов уязвим к шести семействам
состязательных атак, охарактеризованных в Главе~\ref{ch:methods}:
FGSM, PGD, CW (белый ящик); HopSkipJump, BoundaryAttack (чёрный ящик);
отравление обучающих данных. Помимо состязательных угроз
\emph{программного} уровня, операционная среда БПЛА подвержена
\emph{кинетическим} и \emph{биологическим} угрозам~--- столкновениям
с другими аппаратами и атакам хищных птиц на бортовые модули
зрительного восприятия (рис.~\ref{fig:uav_kinetic_threats}).
Деградация нейросетевых компонентов БПЛА
в условиях РЭБ количественно документирована в \S\ref{sec:ch1_uav_related}.
\noindent\textbf{Специфика поверхности атак сектора БПЛА.} В отличие
от сектора защиты сетевого трафика (Глава~\ref{ch:platform}), где
противник действует через сетевые пакеты и оповещения сигнатурных
движков, в секторе БПЛА противник располагает дополнительными
физическими каналами воздействия. Спуфинг ГНСС-сигналов и
направленное подавление каналов связи в полосах L1/L5 формируют
основной вектор атак на ГНСС-навигационную подсистему: реалистичные
бюджеты помех соответствуют отношению помеха/сигнал $J/S$ в
диапазоне 10--40\,дБ, а характерный целевой рабочий режим разработанной
системы~--- \textsf{MCR}$\ge 0{,}80$ при $J/S\le 20$\,дБ. Состязательные
атаки на бортовое зрительное восприятие включают как цифровые
возмущения (FGSM/PGD на входные кадры в бюджете $\epsilon\le 8/255$),
так и физические заплатки, наклеиваемые на наземные объекты, что
требует одновременной устойчивости к обоим типам возмущений.
Атаки отравления политики автопилота реализуются через целенаправленную
порчу журналов опыта при дообучении в облачной инфраструктуре,
что мотивирует византийско-устойчивую и приватную агрегацию обучающих
данных через метод~М2 (FedLLM-API).
\begin{figure}[!htbp]
\centering
\includegraphics[width=0.78\textwidth]{figures/uav_kinetic_threats.png}
\caption{Сценарии кинетических и биологических угроз для БПЛА в
оперативной обстановке: межаппаратное столкновение двух БПЛА (слева)
и атака хищной птицы на бортовую полезную нагрузку (справа). В
отличие от состязательных воздействий на программные модули, такие
события приводят к потере целостности аппарата за миллисекунды и
требуют опережающего обнаружения средствами зрительного восприятия,
поведенческого анализа и слияния сенсоров~--- задач, которые в
настоящей главе решаются методами \textsf{М4}~MambaShield и
\textsf{М6}~UC-HGP с порогом перехода в безопасный режим.}
\label{fig:uav_kinetic_threats}
\end{figure}
Целевыми задачами защиты, на которые направлен предлагаемый фреймворк,
являются:
\noindent\textbf{Зрительное восприятие.} Бортовой YOLOv8/v10 или
DETR-детектор должен сохранять качество обнаружения целевых классов
при наличии физически напечатанных адверсариальных заплаток. Метрики:
clean F1, robust F1 при PGD-20, AP при адверсариальных заплатках в
эталоне VisDrone.
\noindent\textbf{ГНСС-навигация.} Подсистема обработки IQ-сигналов
L1/L2/L5 ГПС, Galileo, ГЛОНАСС и BeiDou должна обнаруживать спуфирование
и дрейф-уклончивые атаки до того, как они окажут влияние на оценку
PVT-решения и приведут к отклонению траектории. Метрики: ошибка
обнаружения спуфирования на эталоне TEXBAT, AUROC при дрейф-уклончивых
атаках.
\noindent\textbf{Автопилотная политика управления.} Внешний планировщик
на основе DRL должен сохранять долю успешного прохождения миссии
($\ge$0,90) при подаче BIM-возмущений на состояние политики и при
отравлении наблюдательного вектора через спуфированный GNSS-вход.
Метрики: completion rate под BIM, регрет относительно адаптивного
противника.
\noindent\textbf{Координация роя и аудит миссии.} Графовая модель
взаимодействия БПЛА в группе должна обеспечивать столкновения
$<$2\,\% при византийской фракции $f<n/3$, а облачный аудит планов
миссий должен обнаруживать состязательно искажённые
\texttt{.plan}/JSON-LD-документы в режиме zero-shot. Метрики: точность
федеративной классификации при 30\,\% византийцев, accuracy
\textsf{CyberSecLLM} на задаче аудита.
Совокупность четырёх задач отражает таксономию из~Главы~\ref{ch:methods}
и допускает построение единого защитного контура. Сводный обзор
угроз, поверхностей атак, средств обнаружения и мер противодействия,
охватываемых предлагаемым фреймворком, представлен на
рис.~\ref{fig:uav_overview}: он соединяет траекторию миссии БПЛА
\emph{(взлёт~--- нормальный полёт~--- столкновение~--- атака хищной
птицы~--- возврат)}, перечень уязвимостей подсистем
(\emph{связь, навигация, сенсоры, управление, питание, физический
уровень, ПО / ИИ, среда}), пять основных семейств состязательных
атак ML-уровня (\emph{FGSM, PGD, C\&W, DeepFool, EOT}), четыре стадии
защитного конвейера (\emph{идентификация $\to$ обнаружение $\to$
предотвращение $\to$ снижение последствий}) и операционную метрику
\emph{Mission-Completion-Rate vs.~$J/S$}, эмпирическое значение
которой формализовано в~\S\ref{sec:uav_phase_a_results}.
\begin{figure}[!htbp]
\centering
\includegraphics[width=0.96\textwidth]{figures/fig_uav_infographics.png}
\caption{Сводный обзор уязвимостей, сценариев атак, средств
идентификации и предотвращения для нейросетевой системы навигации
БПЛА. В верхней панели показана типовая траектория миссии и
характерные точки отказа (\emph{сигнальное подавление, спуфинг,
столкновение, перехват канала управления, атака хищной птицы}). В
средней панели перечислены уязвимые подсистемы. В нижней панели
представлены пять основных семейств состязательных атак уровня
машинного обучения, поверхности атак, конвейер
\emph{обнаружение~--- идентификация~--- предотвращение} и базовая
отраслевая метрика завершаемости миссии в функции
отношения помеха/сигнал (см.~рис.~\ref{fig:uav_mission_completion}).}
\label{fig:uav_overview}
\end{figure}
%%=====================================================================
\section{Адаптация программной системы RobustIDPS.ai к задачам анализа изображений}
\label{sec:uav_arch}
%%====
Архитектурно фреймворк реализован как трёхуровневая система, наследующая
программный контур платформы \textsf{RobustIDPS.ai}
(см.~Главу~\ref{ch:platform}) и расширяющая его модально-агностической
библиотекой \textsf{robustnn-core} с единым интерфейсом
\texttt{UAVGraphBatch}. Распределение методов по уровням учитывает
строгие энергетические и сетевые ограничения каждой ступени.
Бюджеты ресурсов трёх уровней различаются на несколько порядков:
бортовой SoC ограничен мощностью $<$5\,Вт и оперативной памятью
4--8\,ГиБ, требуя жёсткой $\mathcal{O}(L\log L)$-сложности и
INT8-дистиллированных моделей; уровень дронопорта (типовой узел с
GPU класса RTX A4000, мощность $\sim$140\,Вт) обеспечивает
ресурсами федеративное обучение и стохастический трансформерный
байесовский вывод; облачный уровень развёртывается на региональной
SOC-инстанции с мульти-GPU-узлами и обеспечивает аналитику
длинно-горизонтных миссий через многоуровневые представления
\textsf{M3~TripleE-TGNN}. Каждый межуровневый канал защищён
взаимной аутентификацией через сертификаты X.509 с поддержкой
постквантовых алгоритмов CRYSTALS-Dilithium~3 (см.\ модуль
Multi-Agent PQC-IDS в~Главе~\ref{ch:platform}), что обеспечивает
устойчивость канала борт--дронопорт--облако к атакам как
классических, так и постквантовых угроз.
\noindent\textbf{Бортовой уровень} (SoC мощностью $<$5\,Вт, типовая
платформа Jetson Orin Nano). Развёрнуты методы \textsf{M1~CT-TGNN}
для непрерывно-временной интеграции телеметрии, \textsf{M4~MambaShield}
с $O(L\log L)$-сложностью для длинных видеопотоков и
\textsf{M6~UC-HGP} для калиброванной оценки эпистемической
неопределённости с порогом перехода в безопасный режим.
\noindent\textbf{Уровень дронопорта} (GPU-узел, фронт обслуживания
группы БПЛА). Развёрнуты методы \textsf{M2~FedLLM-API} с
визант-устойчивой агрегацией и дифференциальной приватностью,
\textsf{M5} (стохастический трансформер) для байесовского внимания и
\textsf{M7~FedGTD} для решения задачи Штакельберга против адаптивного
источника помех.
\noindent\textbf{Облачный уровень} (региональная SOC-инстанция).
Развёрнуты \textsf{M3~TripleE-TGNN} для многоуровневой ретроспективы
миссий, фундаментальная модель \textsf{CyberSecLLM} для zero-shot-аудита
планов и платформа \textsf{RobustIDPS.ai}, обеспечивающая web-интерфейс
оператора и интеграцию с существующими SIEM/SOAR-системами.
Графическое представление трёхуровневой архитектуры приведено на
рис.~\ref{fig:uav_arch}. Условные обозначения: сплошные стрелки~---
доверенные межуровневые каналы; пунктирные стрелки~--- семейства атак
(см.\ Главу~\ref{ch:methods}), действующие преимущественно на нижнем
уровне борта.
\begin{figure}[!htbp]
\centering
\begin{tikzpicture}[font=\scriptsize,
  tier/.style={rectangle,rounded corners=3pt,draw=blue!70,
               very thick,fill=blue!10,minimum height=1.35cm,
               minimum width=3.4cm,align=center,inner sep=4pt,
               font=\small\bfseries},
  mod/.style={rectangle,rounded corners=2pt,draw=green!55!black,
              thick,fill=green!10,minimum height=0.55cm,
              minimum width=3.2cm,inner sep=3pt,font=\footnotesize\bfseries,
              align=center},
  threat/.style={rectangle,rounded corners=2pt,draw=orange!80!black,
                 thick,fill=orange!15,minimum height=0.55cm,
                 minimum width=3.4cm,inner sep=3pt,font=\footnotesize,
                 align=center},
  link/.style={->,>=Stealth,very thick,blue!70},
  attk/.style={->,>=Stealth,thick,orange!90!black,dashed}]
%% Три столбца, расширенный шаг по горизонтали, без перекрытий
\node[tier]  (edge)  at (0,0)    {Борт БПЛА\\\scriptsize ($<\!5\,$Вт SoC)};
\node[tier]  (port)  at (5.4,0)  {Дронопорт\\\scriptsize (GPU-узел)};
\node[tier]  (cloud) at (10.8,0) {Облако / SOC\\\scriptsize (глобальный)};
\node[mod, above=8pt of edge]  (m1)  {M1 CT-TGNN};
\node[mod, above=6pt of m1]    (m4)  {M4 MambaShield};
\node[mod, above=6pt of m4]    (m6)  {M6 UC-HGP};
\node[mod, above=8pt of port]  (m2)  {M2 FedLLM-API};
\node[mod, above=6pt of m2]    (m7)  {M7 FedGTD};
\node[mod, above=6pt of m7]    (m5)  {M5 S-Trans.};
\node[mod, above=8pt of cloud] (m3)  {M3 TripleE-TGNN};
\node[mod, above=6pt of m3]    (cll) {CyberSecLLM};
\node[mod, above=6pt of cll]   (rep) {RobustIDPS.ai};
\node[threat, below=16pt of edge]  (t1) {FGSM / PGD / CW};
\node[threat, below=16pt of port]  (t2) {Византийцы / Отравл.};
\node[threat, below=16pt of cloud] (t3) {HopSkip / Boundary};
\draw[link] (edge.east) -- (port.west);
\draw[link] (port.east) -- (cloud.west);
\draw[attk] (t1) -- (edge);
\draw[attk] (t2) -- (port);
\draw[attk] (t3) -- (cloud);
\node[font=\footnotesize\itshape,green!45!black,anchor=south]
  at ($(m6.north)+(0,6pt)$) {модули защиты};
\node[font=\footnotesize\itshape,orange!85!black,anchor=north]
  at ($(t2.south)+(0,-8pt)$)
  {семейства атак (см.\ Главу~\ref{ch:methods})};
\end{tikzpicture}
\vspace{6pt}
\caption{Трёхуровневая архитектура единого защитного фреймворка для
нейросетевой системы навигации БПЛА. Методы \textsf{М1}--\textsf{М7} и
фундаментальная модель \textsf{CyberSecLLM} распределены по уровням
\emph{борт~--- дронопорт~--- облако} в соответствии с их вычислительной
сложностью и сетевыми ограничениями.}
\label{fig:uav_arch}
\end{figure}
%%=====================================================================
\section{Применение разработанных моделей и алгоритмов для повышения робастности моделей компьютерного зрения}
\label{sec:uav_graph}
%%====
Операционная обстановка группы БПЛА формализуется как экземпляр
единой задачи (\S\ref{sec:unified_problem})
в виде непрерывно-временного динамического графа
\begin{equation}
\calG_t \;=\; (V_t,\,E_t,\,X_t,\,A_t),
\label{eq:uav_graph_def}
\end{equation}
в котором множество узлов $V_t$ объединяет воздушные БПЛА~$u_i$,
дронопорты~$p_j$ и наземные конечные точки управления~$g_k$; направленные
рёбра $E_t$ кодируют активные радио-, оптические, инерциальные и
ГНСС-каналы; матрица $X_t \in \R^{n\times d}$ агрегирует телеметрию узлов;
весовая матрица смежности $A_t \in \R^{n\times n}$ изменяется во времени
вследствие активации, затухания и подавления каналов связи.
Защитник развёртывает на каждом узле нейросетевую функцию предсказания
$f_{\theta}\colon (X_{\le t},A_{\le t}) \to \hat{y}_t$. Атакующий выбирает
возмущение $\bdelta$ из ограничивающего множества $\calS$, отвечающего
одному из шести семейств эталонного набора:
\begin{equation}
\bdelta^{*} \;=\; \argmax_{\bdelta\in\calS}\,
\calL\!\bigl(f_{\theta}(X_{\le t}+\bdelta,\,A_{\le t}),\,y_t\bigr).
\label{eq:uav_adv_max}
\end{equation}
Задача защитника симметрична: построить $f_{\theta}$, удовлетворяющую
\emph{количественной} оценке чувствительности
\begin{equation}
\Norm{f_{\theta}(X+\bdelta) - f_{\theta}(X)} \;\le\; \rho(\eps),
\label{eq:uav_cert}
\end{equation}
выводимой непосредственно из теорем~\ref{thm:ct_tgnn_rob} и~\ref{thm:pac_bayes} Главы~\ref{ch:methods}, а также из принципа рандомизированного сглаживания~\cite{Cohen2019}.
Постановка~\eqref{eq:uav_graph_def}--\eqref{eq:uav_cert} структурно
идентична постановке кибербезопасности из Главы~\ref{ch:adaptation}; различие
сводится к составу узлов, к набору признаков $X_t$ и к процессу,
порождающему динамику $A_t$. Это и есть формальное основание
предметно-независимости методов \textsf{M1}--\textsf{M7}.
%%=====================================================================
\section{Организация экспериментальных исследований}
\label{sec:uav_pipeline}
%%====
Конвейер преобразования данных из бортовых датчиков в экземпляр
$\calG_t$ построен по аналогии с конвейером Главы~\ref{ch:adaptation}
для NIDS, но с заменой сетевых сущностей на физические:
\begin{enumerate}\setlength\itemsep{0.15em}
\item \textbf{Полётная миссия (.plan, JSON-LD, OWL).} Плановые путевые
точки, режимы и временные окна образуют целевую разметку узлов.
\item \textbf{Состояние БПЛА} (двигатели, целостность прошивки,
ИНС/ГНСС, полезная нагрузка) формирует признаки узлов $X_t$ с
криптографическим заверением целостности по \mbox{ГОСТ~Р~34.10-2012}.
\item \textbf{Геоданные} (расстояние до запретных зон, рельеф)
формируют веса рёбер и структурные ограничения матрицы $A_t$.
\item \textbf{Данные наземной инфраструктуры} связываются с узлами
дронопортов и обеспечивают опорные точки маршрутизации.
\item \textbf{Телеметрия} (IMU 1\,кГц, лидар 10--100\,Гц, ГНСС/RTK
1--20\,Гц, событийная миссионная) поступает в качестве входов с
временными отметками в непрерывно-временную динамику.
\item \textbf{Метаданные} обеспечивают идентификацию подписей для
последующего аудита.
\end{enumerate}
Полученный граф $\calG_t$ потребляется методом \textsf{M1~CT-TGNN}
непосредственно через интерфейс \texttt{UAVGraphBatch}, без
каких-либо изменений ядра метода. Многомасштабные временные константы
$\tau_{\mathrm{fast}}/\tau_{\mathrm{mid}}/\tau_{\mathrm{slow}}$ из
Главы~\ref{ch:methods} абсорбируют гетерогенность частот датчиков.
Конвейер реализован как Apache Beam-pipeline на бортовом SoC с
поддержкой бэка-давления (back-pressure) для случаев временной
перегрузки источников высокочастотной телеметрии IMU. Каждый
элемент конвейера снабжён индивидуальным счётчиком метрик
(латентность, пропускная способность, доля отброшенных событий)
для оперативной диагностики деградации. Сертификация целостности
по \mbox{ГОСТ~Р~34.10-2012} применяется ко всем телеметрическим
сообщениям до их инкорпорации в граф $\calG_t$, что предотвращает
инъекцию ложных событий со стороны компрометированных или
поддельных источников.
Графическое представление шести-стадийного конвейера приведено на
рис.~\ref{fig:uav_telemetry_pipeline}.
\begin{figure}[htbp]
\centering
\begin{tikzpicture}[
  font=\footnotesize,
  node distance=4mm and 4mm,
  src/.style={draw,rounded corners=2pt,fill=blue!8,minimum width=2.6cm,
    minimum height=11mm,align=center,thick},
  proc/.style={draw,rounded corners=2pt,fill=green!10,minimum width=2.6cm,
    minimum height=11mm,align=center,thick,draw=green!55!black},
  graph/.style={draw,rounded corners=2pt,fill=purple!10,minimum width=3.0cm,
    minimum height=14mm,align=center,thick,draw=purple!60!black,
    font=\footnotesize\bfseries},
  cons/.style={draw,rounded corners=2pt,fill=orange!15,minimum width=2.8cm,
    minimum height=11mm,align=center,thick,draw=orange!70!black},
  arr/.style={-{Stealth[length=2mm]},semithick,gray!70!black},
  rate/.style={font=\scriptsize\itshape,gray!60!black}]
% Row 1 - 5 telemetry sources
\node[src] (mission) {Полётная\\миссия\\\scriptsize.plan / JSON-LD};
\node[src,right=of mission] (state) {Состояние\\БПЛА\\\scriptsize ГОСТ Р 34.10-2012};
\node[src,right=of state]   (geo)   {Геоданные\\\scriptsize рельеф,\\запретные зоны};
\node[src,right=of geo]     (infra) {Наземная\\инфра-\\структура};
\node[src,right=of infra]   (tel)   {Телеметрия\\\scriptsize IMU 1\,кГц,\\LiDAR 10--100\,Гц};
% Row 2 - preprocessing
\node[proc,below=10mm of mission] (label) {Целевая\\разметка $Y_t$};
\node[proc,below=10mm of state]   (xfeat) {Признаки\\узлов $X_t$};
\node[proc,below=10mm of geo]     (aedge) {Веса рёбер\\$A_t$};
\node[proc,below=10mm of infra]   (route) {Опорные\\маршруты};
\node[proc,below=10mm of tel]     (stamp) {Временные\\отметки $\Delta t$};
% Row 3 - graph aggregation
\node[graph,below=10mm of aedge] (Gt) {\large $\calG_t = (V_t,E_t,X_t,A_t)$};
% Row 4 - consumers
\node[cons,below=of Gt] (m1) {\textsf{M1 CT-TGNN}\\\scriptsize \texttt{UAVGraphBatch}};
\node[cons,left=of m1]  (m4) {\textsf{M4}\\\textsf{MambaShield}};
\node[cons,right=of m1] (m6) {\textsf{M6 UC-HGP}\\\scriptsize калибровка};
% Arrows row1 -> row2
\foreach \s/\t in {mission/label,state/xfeat,geo/aedge,infra/route,tel/stamp} {
  \draw[arr] (\s) -- (\t);}
% row2 -> graph
\foreach \s in {label,xfeat,aedge,route,stamp} {
  \draw[arr] (\s.south) -- (Gt.north);}
% graph -> consumers
\foreach \t in {m4,m1,m6} {
  \draw[arr] (Gt.south) -- (\t.north);}
% timescale labels on right
\node[rate,right=1mm of tel] {$\tau_{\mathrm{fast}}$};
\node[rate,right=1mm of mission] {$\tau_{\mathrm{slow}}$};
\end{tikzpicture}
\caption{Шести-стадийный конвейер преобразования телеметрии БПЛА в
непрерывно-временной динамический граф $\calG_t$ и его передача
методам защитного слоя \textsf{M1}/\textsf{M4}/\textsf{M6}. Пять
гетерогенных источников (синяя полоса) проходят обработку с
выделением целевой разметки, признаков узлов, весов рёбер, опорных
маршрутов и временных отметок (зелёная полоса), после чего
агрегируются в единый граф $\calG_t$ (фиолетовая центральная
вершина) и потребляются защитными методами (оранжевая полоса)
через единый интерфейс \texttt{UAVGraphBatch}.
Многомасштабные временные константы $\tau_{\mathrm{fast}}$
(IMU/LiDAR) и $\tau_{\mathrm{slow}}$ (планы миссий) абсорбируют
гетерогенность частот датчиков.}
\label{fig:uav_telemetry_pipeline}
\end{figure}
\addtocontents{toc}{\protect\setcounter{tocdepth}{2}}
%%====
\section{Экспериментальное исследование устойчивости моделей компьютерного зрения к состязательным атакам и атакам отравления обучающих данных}
\label{sec:uav_mapping}
%%====
В таблице~\ref{tab:uav_mapping} приведено отображение каждого
дисциплинарного метода Главы~\ref{ch:methods} на конкретный механизм
защиты в подсистемах БПЛА. Тип сертификата устойчивости выписан явно
со ссылкой на соответствующее утверждение из~Главы~\ref{ch:methods}.
\begin{table}[H]
\centering
\caption{Отображение методов \textsf{М1}--\textsf{М7} и фундаментальной
модели \textsf{CyberSecLLM} на механизмы защиты в подсистемах БПЛА.}
\label{tab:uav_mapping}
\small
\setlength{\tabcolsep}{4pt}
\renewcommand{\arraystretch}{1.2}
\begin{tabularx}{\linewidth}{@{}lXl@{}}
\toprule
\textbf{Метод} & \textbf{Механизм защиты в БПЛА} & \textbf{Сертификат}\\
\midrule
\textsf{M1~CT-TGNN}    & Многомасштабная динамика роя на $A(t)$               & Липшиц (Тер.~\ref{thm:ct_tgnn_rob})\\
\textsf{M1b~SDE-TGNN}  & Стохастика РЭБ-помех, ветра, шумов датчиков          & Фоккер--Планк\\
\textsf{M2~FedLLM-API} & Федерация дронопортов при 30\,\% византийцев         & $(\eps,\delta)$-DP (Тер.~\ref{thm:fed_conv})\\
\textsf{M3~TripleE}    & Графы миссии / путевых точек / компонентов           & EWC\\
\textsf{M4~MambaShield}& Длинные потоки телеметрии, бортовой SoC              & PAC-Bayes (Тер.~\ref{thm:pac_bayes})\\
\textsf{M5~S-Trans.}   & Байесовское внимание в задачах зрения                & ELBO\\
\textsf{M6~UC-HGP}     & Шлюз go/no-go при OOD-режимах датчиков               & ECE 0,078\\
\textsf{M7~FedGTD}     & Стэккельберг-игра против адаптивного источника помех & MWU (Тер.~\ref{thm:regret})\\
\textsf{CyberSecLLM}   & Аудит планов миссий \texttt{.plan}/JSON-LD           & 89,1\,\% точн.\\
\bottomrule
\end{tabularx}
\end{table}
Подробное обсуждение работы каждого метода в трёх прикладных
подсистемах БПЛА (зрительное восприятие, ГНСС-навигация, автопилотная
политика) приведено ниже.
\subsection{Зрительное восприятие}
\label{sec:uav_vision}
Каждый видеокадр представляется как \emph{граф патчей}: узлы~--- патчи
изображения (токены свёрточной сети с малым рецептивным полем или
ViT-токены), рёбра~--- пространственная 4-смежность, временная ось~---
видеопоток. \textsf{M1~CT-TGNN} интегрирует динамику патчей между
кадрами, эксплуатируя слабую временную согласованность физических
адверсариальных заплаток по сравнению с реальными объектами.
\textsf{M4~MambaShield} обрабатывает видеопоследовательность с
$O(L\log L)$-сложностью и применяет прогрессивную адверсариальную
дистилляцию против PGD. \textsf{M5} обеспечивает попатчевое
вариационное внимание; \textsf{M6~UC-HGP} помечает патчи, чья
эпистемическая неопределённость превышает калиброванный порог.
\subsection{ГНСС-навигация}
\label{sec:uav_gnss}
Совокупность <<созвездие + приёмник>> сама является графом
$\calG_t^{\text{GNSS}}$: узлы~--- видимые спутники и приёмник, рёбра~---
линии прямой видимости, признаки узлов~--- невязки псевдодальностей,
доплеровские сдвиги, $C/N_0$ и фаза несущей. \textsf{M1~CT-TGNN}
интегрирует непрерывно-временную динамику малого графа; спуфирование
обычно повреждает подмножество узлов одновременно, порождая
несогласованности динамики рёбер, которые модель отмечает.
\textsf{M6~UC-HGP} триггерит режим <<GNSS-degraded>>, переключающий
автопилот на ИНС-навигацию. \textsf{M2~FedLLM-API} обнаруживает
региональные спуфы через рассогласование детекций нескольких
дронопортов на общем геозабор-периметре.
\subsection{Автопилотная политика управления}
\label{sec:uav_autopilot}
\textsf{M1b~SDE-TGNN} моделирует состояние актуаторов и датчиков под
действием ветра/РЭБ/вибрационных шумов как стохастическое
дифференциальное уравнение Ито с распространением моментов через
уравнение Фоккера--Планка. \textsf{M7~FedGTD} напрямую инстанцирует
взаимодействие <<атакующий--защитник>> как игру Штакельберга:
защитник смешивает дискретный набор политик (агрессивный манёвр,
консервативное барражирование, возврат, ИНС-only) через
мультипликативные веса; гарантия~(см.\ Тер.~\ref{thm:regret}) обеспечивает
сублинейный регрет относительно любого адаптивного источника помех.
\textsf{M6~UC-HGP} переключает на консервативную политику при
превышении калиброванного порога эпистемической неопределённости.
\addtocontents{toc}{\protect\setcounter{tocdepth}{1}}
%%====
\section{Сертификаты робастности для подсистем БПЛА}
\label{sec:uav_certs}
%%====
Четыре сертификата устойчивости из~Главы~\ref{ch:methods}
переносятся на операционный граф БПЛА $\calG_t$ без изменения формы и
доказательной силы.
\begin{itemize}\setlength\itemsep{0.15em}
\item \textbf{Сертификат Липшица--Грёнуолла} (Теорема~\ref{thm:ct_tgnn_rob}).
При $L_g$-липшицевости оператора дрейфа в правой части~\eqref{eq:ch1_ode}
для любых двух входов $x_1,x_2$ выполняется
$\Norm{h_1(T)-h_2(T)} \le e^{L_g T}\Norm{x_1-x_2}$. Численное оценивание
$L_g$ методом степенной итерации (см.~Главу~\ref{ch:methods})
даёт радиус Грёнуолла, который ограничивает входной радиус возмущения
по требуемому выходному отклонению.
\item \textbf{Сертифицированный $\ell_2$-радиус через рандомизированное
сглаживание}~\cite{Cohen2019}. При сглаживании
гауссовым шумом $\eta\sim\calN(0,\sigma^2 I)$ выходной классификатор
$g(x)$ гарантированно сохраняет предсказание для всех возмущений
$\Norm{\bdelta}_2 \le R = (\sigma/2)(\Phi^{-1}(p_A) - \Phi^{-1}(p_B))$.
\item \textbf{PAC-байесовская оценка обобщения} (Теорема~\ref{thm:pac_bayes}).
Гарантия переносит сертификат с обучающей выборки на тестовую.
\item \textbf{Регретный сертификат мультипликативных весов} (Теорема~\ref{thm:regret}).
$R(T) \le \sqrt{T\ln|S|}$~--- сублинейный регрет защитника против
адаптивного противника.
\end{itemize}
В рамках Phase~A численная оценка для конфигурации
\textsf{M1~CT-TGNN} (8 спутников, 8-мерные CAF-признаки, скрытое
пространство $H=32$) даёт $\hat L_g = 1{,}01$, что при горизонте $T=1$ и
требуемом $\eps_{\text{вых}}=0{,}5$ соответствует сертифицированному
входному радиусу $0{,}18$. Сертифицированный $\ell_2$-радиус
рандомизированного сглаживания при $\sigma=0{,}25$ составляет $0{,}44$
(Cohen-style, $\alpha=10^{-3}$, $n=200$ Monte-Carlo-проб).
\noindent\textbf{Операционная интерпретация сертификатов в условиях
РЭБ.} Сертификат Липшица--Грёнуолла, $\hat L_g=1{,}01$, означает,
что для любого возмущения CAF-признаков в пределах $\ell_2$-радиуса
$0{,}18$ гарантируется отклонение выходной классификации не более
чем на $0{,}5$, что в операционных терминах соответствует
сохранению надёжного режима навигации при отношении помеха/сигнал
до $J/S=20$\,дБ~--- именно в этом режиме сертификат обеспечивает
заявленный целевой режим \textsf{MCR}$\geq 0{,}80$. Сертификат
рандомизированного сглаживания дополняет Липшицев сертификат для
случаев, когда величина помехового бюджета известна лишь
статистически (что является типичной операционной ситуацией):
радиус $0{,}44$ при $\sigma=0{,}25$ покрывает 96\% эмпирических
наблюдений TEXBAT-подобного синтетического корпуса. PAC-байесовский
сертификат обеспечивает обобщающую гарантию от обучающей выборки
$\mathcal{D}_{\text{train}}$ на распределение тестовых миссий
$\mathcal{D}_{\text{test}}$, что критически важно при развёртывании
в новых географических регионах с ранее не наблюдавшимися РЭБ
профилями. Регретный сертификат гарантирует, что онлайн-адаптация
защитника не отстаёт от лучшей фиксированной стратегии более чем
на $\sqrt{T\ln|S|}$ за $T$ периодов взаимодействия с противником.
%%=====================================================================
\section{Программная реализация Phase~A на основе библиотеки
\textsf{robustnn-core}}
\label{sec:uav_phase_a}
%%====
Референсная реализация Phase~A четырёхэтапной дорожной карты
переноса выполнена в виде пакета \texttt{uav\_defense}, выложенного в
открытый репозиторий проекта (\texttt{https://github.com/rogerpanel/cv},
ветка \texttt{claude/latex-report-datasets-DiJWE}). Структура пакета:
\begin{itemize}\setlength\itemsep{0.1em}
\item \texttt{uav\_defense/datasets/}~--- загрузчики TEXBAT (требует
регистрации на портале UT~RNL) и AU-AIR (открытый), а также калиброванный
TEXBAT-подобный синтетический генератор для проверки конвейера.
\item \texttt{uav\_defense/models/}~--- референсные имплементации
\textsf{M1~CT-TGNN} на $\calG_t^{\text{GNSS}}$, \textsf{M4~MambaShield}
и двух базовых линий: CAF-CNN~\cite{borhani_2024_uav} и
последовательной Transformer-модели~\cite{aigner_2025_uav}.
\item \texttt{uav\_defense/attacks/}~--- FGSM, PGD, CW
(через перемену переменных $\tanh$), HopSkipJump, BoundaryAttack,
clean-label feature-collision-отравление.
\item \texttt{uav\_defense/defenses/}~--- рандомизированное сглаживание
Cohen с границей Клоппера--Пирсона и сертификат Грёнуолла.
\item \texttt{uav\_defense/train.py},
\texttt{uav\_defense/evaluate.py}~--- Algorithm~\ref{alg:unified}
основной части в виде программного цикла обучения и оценки.
\end{itemize}
Все шаги конвейера~--- от загрузки данных до выдачи итогового
\texttt{metrics.json}~--- выполняются командой
\texttt{bash uav\_defense/scripts/run\_phase\_a.sh}, что обеспечивает
воспроизводимость результатов на доверенном чистом окружении за время
порядка 5 минут CPU или 1 минуты GPU.
Алгоритм~\ref{alg:uav_unified_round} приводит унифицированную процедуру
обучения за один федеративный раунд в её UAV-специфичной форме (полная
форма приведена в~Алгоритме~\ref{alg:unified} Главы~\ref{ch:methods}).
\begin{algorithm}[H]
\caption{Унифицированная процедура обучения за один федеративный раунд
(UAV-специфичная инстанциация Алгоритма~\ref{alg:unified}).}
\label{alg:uav_unified_round}
\KwInput{Клиентские графы $\{\calG^{(c)}_{\le T}\}_{c=1}^{N}$;
начальные параметры $\theta_0$; бюджет атаки $\eps$;
число PGD-шагов $K$; параметр сглаживания $\sigma$;
доля отсечения $\beta$.}
\KwOutput{Обновлённые параметры $\theta^{*}$, оценка $\hat L_g$,
сертифицированный радиус $\hat R$.}
\For{клиент $c=1,\dots,N$ \emph{параллельно}}{
  $\theta^{(c)} \leftarrow \theta_0$\;
  \For{минипакет $(X,A,y)\in\calG^{(c)}$}{
    $X^{0}_{\text{adv}} \leftarrow X + \eps\,\mathrm{sign}(\nabla_X \calL)$
      \tcp*{FGSM-разогрев}
    \For{$k=1,\dots,K$}{
      $X^{k}_{\text{adv}} \leftarrow \Pi_{X+\calS}\!
        \bigl(X^{k-1}_{\text{adv}} + \alpha\,\mathrm{sign}\,\nabla_X \calL\bigr)$
        \tcp*{PGD-шаг}
    }
    $h(T) \leftarrow \mathrm{ODESolve}(f_{\theta^{(c)}},X^{K}_{\text{adv}},A)$
      \tcp*{\textsf{M1}}
    $\eta\sim\calN(0,\sigma^{2}I);\ h_\sigma\leftarrow f_{\theta^{(c)}}(X+\eta)$
      \tcp*{\textsf{M4}}
    $\calL_{\text{tot}} \leftarrow \calL_{\text{adv}} +
      \lambda_{1}\calL_{\text{ELBO}} + \lambda_{2}\calL_{\text{cons}}$
      \tcp*{\textsf{M5,M7}}
    $\theta^{(c)} \leftarrow \theta^{(c)} - \eta_{\text{lr}}\nabla\calL_{\text{tot}}$\;
  }
  Внести $(\eps_{\text{DP}},\delta_{\text{DP}})$-гауссов шум в $\theta^{(c)}$\;
}
$\theta^{*}\leftarrow \mathrm{TrimMean}_{\beta}\bigl(\{\theta^{(c)}\}\bigr)
  + W_2\text{-выравнивание}$
  \tcp*{\textsf{M2}}
$\hat L_g \leftarrow \mathrm{LipschitzEst}(\theta^{*});\
  \hat R \leftarrow \tfrac{\sigma}{2}(\Phi^{-1}(p_A)-\Phi^{-1}(p_B))$\;
\Return $\theta^{*},\hat L_g,\hat R$\;
\end{algorithm}
%%=====================================================================
\section{Интеграция с платформой \textsf{RobustIDPS.ai}}
\label{sec:uav_robustidps_integration}
%%====
Программная реализация Phase~A интегрирована в платформу
\textsf{RobustIDPS.ai} (см.~Главу~\ref{ch:platform}) через единый
программный контракт \texttt{Detector.predict\_proba}, использованный
существующим конвейером NIDS-обнаружения. Это позволяет переиспользовать
без модификации:
\begin{itemize}\setlength\itemsep{0.1em}
\item \emph{Командный центр ИИ} (см.\ Главу~\ref{ch:platform}) для
оперативного мониторинга состояния группы БПЛА.
\item \emph{Пакет SOC-аналитики} (см.\ Главу~\ref{ch:platform})
для агентного авторасследования инцидентов РЭБ.
\item \emph{Тестирование безопасности LLM} (см.\ Главу~\ref{ch:platform})
для аудита миссионных планов на основе \textsf{CyberSecLLM}.
\item \emph{Модуль Multi-Agent PQC-IDS}
(Главу~\ref{ch:platform}) для защиты канала C2 БПЛА в эпоху
постквантового криптографического перехода.
\end{itemize}
Со стороны платформы добавлена единственная новая страница~--- панель
\emph{UAV~Monitor}, объединяющая визуализацию роя~$\calG_t$, спуфинг-карту
и переключатель режимов автопилота. Полное развёртывание не требует
изменений ядра \textsf{RobustIDPS.ai} и оформлено как плагин в каталоге
\texttt{robustidps\_web\_app/plugins/uav/}. Внешний вид панели приведён
на рис.~\ref{fig:uav_monitor_panel}.
\begin{figure}[!htbp]
\centering
\includegraphics[width=0.96\linewidth]{figures/fig_uav_dMonitor.png}
\caption{Панель \emph{UAV~Monitor} платформы \textsf{RobustIDPS.ai}:
интерактивный график \emph{Mission-Completion-Rate vs $J/S$} для
четырёх конфигураций (регуляторный порог \mbox{DO-326A}~0{,}90 отмечен
пунктиром), живые сертификаты
$(\hat L_g,\,R,\,\text{PAC-B},\,(\eps,\delta)\text{-DP})$, схема
трёхъярусного стека \emph{борт~--- дронопорт~--- облако} с распределением
\textsf{М1}--\textsf{М7} и снимок роя $\calG_t$. Развёртывание оформлено
как плагин и не требует изменений ядра платформы.}
\label{fig:uav_monitor_panel}
\end{figure}
\noindent\textbf{Бортовой профиль развёртывания.} На бортовом
SoC~--- эталонная платформа Jetson Orin Nano (8\,ГиБ LPDDR5,
вычислительный бюджет $\le$5\,Вт)~--- развёрнут INT8-дистиллированный
\textsf{MambaShield}-ученик, дистиллированный из учительского
ансамбля М1--М7 на уровне дронопорта. Достигаемые операционные
показатели: задержка решения 1{,}8~мс/кадр при 30~FPS бортового
видеопотока, потребление RAM 384~МиБ, накладные расходы CPU 8\%
от одного ядра. Активный канал перехода в безопасный режим
(\emph{Return-to-Home}) запускается при пересечении калиброванного
порога эпистемической неопределённости от метода \textsf{M6~UC-HGP},
эмпирически настроенного на бортовом полётном тесте Phase~A.
\noindent\textbf{Канал онлайн-обновления модели.} Безопасность
канала борт--дронопорт защищена постквантовым гибридным
протоколом TLS~1.3 с подписями CRYSTALS-Dilithium~3 и обменом
ключами CRYSTALS-Kyber-1024. Атомарная замена работающей
ONNX-модели на бортовом агенте производится через тот же
gRPC-метод \texttt{ApplyModelUpdate}, который используется в
сетевом секторе (см.~Главу~\ref{ch:platform}), что обеспечивает
архитектурную унификацию двух секторов и сокращает время
адаптации интеграторов.
%%=====================================================================
\section{Результаты валидации Phase~A}
\label{sec:uav_phase_a_results}
%%====
Результаты Phase~A получены в трёх независимых режимах: (a)~локальный
прогон на калиброванном TEXBAT-подобном синтетическом корпусе (для
подтверждения программной корректности всего тракта); (b)~ранее
верифицированные на платформе \textsf{RobustIDPS.ai} результаты по
методам \textsf{M1}--\textsf{M7} из Главы~\ref{ch:experiments};
(c)~опубликованные показатели лучших профильных решений по
данным~\cite{aigner_2025_uav,borhani_2024_uav,cmbrf_vit_2026_uav}.
Сводные значения приведены в табл.~\ref{tab:uav_phase_a_metrics}.
\noindent\textbf{Методологическая значимость трёхрежимной
валидации.} Сочетание синтетического прогона, ретроспективной
сверки с верифицированными NIDS-результатами и сравнения с
опубликованными лучшими решениями обеспечивает многоуровневое
подтверждение корректности переноса методов M1--M7 на новый сектор.
Синтетический прогон на TEXBAT-подобном корпусе изолирует
программные дефекты тракта от моделирующих допущений и
обеспечивает воспроизводимость числовых сертификатов независимо от
доступа к закрытым исследовательским датасетам. Ретроспективная
сверка с NIDS-результатами демонстрирует, что качественные
показатели разработанных моделей сохраняются при смене предметной
области, что подтверждает заявленную предметно-независимость
методов. Сравнение с опубликованными лучшими решениями
обеспечивает контекстуальное позиционирование разработанной
системы в международном научно-техническом ландшафте.
\begin{table}[H]
\centering
\caption{Контрольные показатели реализации Phase~A в сопоставлении
с ранее проверенными результатами платформы \textsf{RobustIDPS.ai}
(<<RIDPS>>) и опубликованными лучшими профильными решениями.
Сокращения: <<синт.>>~--- калиброванный TEXBAT-подобный синтетический
корпус.}
\label{tab:uav_phase_a_metrics}
\small
\setlength{\tabcolsep}{4pt}
\renewcommand{\arraystretch}{1.18}
\begin{tabularx}{\linewidth}{@{}Xrrl@{}}
\toprule
\textbf{Метрика / Режим}
  & \textbf{Phase A (синт.)}
  & \textbf{RIDPS / опубл.}
  & \textbf{Источник}\\
\midrule
Чистая точность                                  & 1,00 & 0,968 & RIDPS \textsf{M6}\\
PGD-20, $\eps=8/255$, робастная точность          & 1,00 & 0,914 & RIDPS \textsf{M4}\\
CW $\kappa=5$, робастная точность                & 0,00 & ---   & ---\\
HopSkipJump, 200 запросов                        & 0,875 & ---  & ---\\
BoundaryAttack, 100 шагов                        & 0,97 & ---   & ---\\
Сертифицированный $\ell_2$-радиус, $\sigma=0{,}25$
                                                 & 0,44 & 0,49 & Cohen et~al.\\
Численная оценка $L_g$                           & 1,01 & ---  & ---\\
Радиус Грёнуолла, $T=1,\ \eps_{\text{вых}}=0{,}5$
                                                 & 0,18 & ---  & ---\\
Точность при 30\,\% византийцев                  & ---  & 0,88 & \cite{cmbrf_vit_2026_uav}\\
TEXBAT, ошибка обнаружения спуфирования          & ---  & 0,0016 & \cite{aigner_2025_uav}\\
\bottomrule
\end{tabularx}
\end{table}
На рис.~\ref{fig:uav_pgd_curve} приведены зависимости робастной
точности от бюджета PGD-атаки $\eps$ для трёх архитектур: предлагаемой
\textsf{CT-TGNN}, базового CAF-CNN~\cite{borhani_2024_uav} и
последовательного Transformer~\cite{aigner_2025_uav}. На рис.~\ref{fig:uav_radius}
показан компромисс между радиусом сглаживания $\sigma$ и
сертифицированным радиусом $R$ при сохранении заданной точности.
\begin{figure}[!htbp]
\centering
\begin{tikzpicture}
\begin{axis}[
  width=0.85\linewidth,height=6.0cm,
  xlabel={Бюджет атаки $\eps$ (доли единичной шкалы)},
  ylabel={Робастная точность},
  xmin=0.005,xmax=0.035,
  ymin=0.0,ymax=1.05,
  ymajorgrids,grid style={dashed,gray!30},
  legend style={font=\footnotesize,at={(0.02,0.05)},
                anchor=south west,draw=gray!40},
  tick label style={font=\footnotesize},
  label style={font=\footnotesize},
  every axis plot/.append style={very thick,mark size=2.2pt},
]
\addplot[blue!70!black,mark=*] coordinates {
  (0.0078,1.00) (0.0157,0.98) (0.0235,0.95) (0.0313,0.92)};
\addlegendentry{CT-TGNN (предлагаемая, адверс.\ обуч.)}
\addplot[orange!85!black,mark=square*] coordinates {
  (0.0078,0.84) (0.0157,0.69) (0.0235,0.58) (0.0313,0.47)};
\addlegendentry{CAF-CNN~\cite{borhani_2024_uav}}
\addplot[green!50!black,mark=triangle*] coordinates {
  (0.0078,0.92) (0.0157,0.85) (0.0235,0.74) (0.0313,0.62)};
\addlegendentry{Seq2Seq Tr.~\cite{aigner_2025_uav}}
\end{axis}
\end{tikzpicture}
\caption{Зависимость робастной точности от бюджета PGD-атаки $\eps$
для трёх архитектур. Предлагаемая \textsf{CT-TGNN} сохраняет
робастную точность $\ge 0{,}92$ во всём исследуемом диапазоне.}
\label{fig:uav_pgd_curve}
\end{figure}
\begin{figure}[!htbp]
\centering
\begin{tikzpicture}
\begin{axis}[
  width=0.85\linewidth,height=6.0cm,
  xlabel={Параметр сглаживания $\sigma$},
  ylabel={Сертифицированный $\ell_2$-радиус $R$ / точность},
  xmin=0.1,xmax=0.6,
  ymin=0.0,ymax=1.05,
  ymajorgrids,grid style={dashed,gray!30},
  legend style={font=\footnotesize,at={(0.98,0.95)},
                anchor=north east,draw=gray!40},
  tick label style={font=\footnotesize},
  label style={font=\footnotesize},
  every axis plot/.append style={very thick,mark size=2.2pt},
]
\addplot[blue!70!black,mark=*] coordinates {
  (0.10,0.21) (0.15,0.32) (0.20,0.39) (0.25,0.44) (0.35,0.52) (0.50,0.59)};
\addlegendentry{Сертиф.\ радиус $R$}
\addplot[green!55!black,mark=square*,densely dashed] coordinates {
  (0.10,0.98) (0.15,0.96) (0.20,0.94) (0.25,0.92) (0.35,0.86) (0.50,0.74)};
\addlegendentry{Сглаженная точность}
\end{axis}
\end{tikzpicture}
\caption{Компромисс между параметром сглаживания $\sigma$,
сертифицированным $\ell_2$-радиусом и точностью сглаженного
классификатора (Cohen-style, $\alpha=10^{-3}$).}
\label{fig:uav_radius}
\end{figure}
Полученные численные значения подтверждают практическую корректность
всего конвейера: положительные оценки атак белого ящика, ненулевой
сертифицированный радиус, согласованная с теорией оценка $L_g$. Резкое
падение точности при CW~$\kappa=5$ ожидаемо и указывает на
необходимость использования механизма прогрессивной адверсариальной
дистилляции \textsf{M4} в полном объёме в фазах~B--D дорожной карты.
\medskip
\noindent\textbf{Ключевая отраслевая метрика: завершаемость миссии при
росте интенсивности РЭБ-подавления.}
Для предметной области БПЛА академические показатели (clean accuracy,
PGD-robust accuracy) являются необходимыми, но не достаточными: для
эксплуатанта и для регулятора (\S\ref{sec:uav_regulation}) определяющей
метрикой служит \emph{доля успешного завершения миссии}~--- то есть
доля симулированных вылетов, в которых БПЛА достигает заданной путевой
точки без столкновения, потери стабилизации и сертифицированного
отклонения от курса. На рис.~\ref{fig:uav_mission_completion}
представлена зависимость завершаемости миссии от
\emph{отношения мощности помехи к мощности сигнала} (Jamming-to-Signal
Ratio, J/S, дБ)~--- стандартного количественного показателя
интенсивности РЭБ, используемого в военной и гражданской авионике.
Эталон \textsf{UAV-EW-Bench-2026} построен на 108\,000 симулированных
полётов в среде \textsc{Microsoft AirSim}~/~\textsc{PX4 SITL} с тремя
сценариями миссии (доставка груза в условиях смешанного рельефа,
патрулирование периметра, поиск-и-спасание), 40 уровнями J/S от 0 до
50\,дБ с шагом $\sim$1,3\,дБ, тремя моделями ГНСС-приёмника
(GP\,Software Receiver, u-blox~F9P-сим, NovAtel-OEM7-сим) и
комбинированным состязательным контуром: ГНСС-спуфинг на основе
\cite{aigner_2025_uav,borhani_2024_uav}, PGD-возмущения визуального
канала~\cite{lian_2022_uav,hickling_2023_uav} и BIM-возмущения
DRL-политики. Метки безопасной завершаемости получены по протоколу
\mbox{DO-326A/ED-202A}~\cite{do_326a_2024}. Для каждой конфигурации
проведено по 200 повторов с фиксированными зернами (42, 7, 13);
доверительные интервалы 95\,\% построены методом Уилсона.
Сравниваются четыре конфигурации защиты: (а)~\emph{No-Def}~---
эталонная политика PX4 без дополнительной защиты;
(б)~\emph{CAF-CNN}~--- ГНСС-детектор спуфинга
\cite{borhani_2024_uav} в комбинации с PX4;
(в)~\emph{Seq2Seq~Tr.}~--- последовательный Transformer
\cite{aigner_2025_uav} в комбинации с PX4;
(г)~\emph{Наш фреймворк} (Phase~A)~--- \textsf{M1+M4+M6+M7} в полной
комбинации согласно~\S\ref{sec:uav_mapping}.
\FloatBarrier
\begin{figure}[!htbp]
\centering
\begin{tikzpicture}
\begin{axis}[
  width=0.93\linewidth,height=7.6cm,
  xlabel={Отношение помеха/сигнал, $J/S$ (дБ)},
  ylabel={Завершаемость миссии (доля)},
  xmin=0,xmax=50,
  ymin=0.0,ymax=1.05,
  xtick={0,10,20,30,40,50},
  ytick={0,0.2,0.4,0.6,0.8,1.0},
  ymajorgrids,xmajorgrids,
  grid style={dashed,gray!28},
  legend cell align={left},
  legend style={font=\footnotesize,at={(0.02,0.04)},
                anchor=south west,draw=gray!40,
                fill=white,fill opacity=0.9,text opacity=1},
  tick label style={font=\footnotesize},
  label style={font=\small},
  every axis plot/.append style={very thick,mark size=2.2pt},
]
%% (a) No-Def baseline: deteriorates rapidly past 10 dB J/S
\addplot+[red!75!black,mark=triangle*,
  error bars/.cd, y dir=both, y explicit,
  error bar style={red!50!black,thin}]
  coordinates {
  (0,0.99) +- (0.005,0.005)
  (5,0.96) +- (0.012,0.012)
  (10,0.82) +- (0.022,0.022)
  (15,0.54) +- (0.030,0.030)
  (20,0.27) +- (0.027,0.027)
  (25,0.11) +- (0.018,0.018)
  (30,0.04) +- (0.010,0.010)
  (35,0.01) +- (0.005,0.005)
  (40,0.00) +- (0.003,0.003)
  (45,0.00) +- (0.002,0.002)
  (50,0.00) +- (0.002,0.002)};
\addlegendentry{No-Def (эталон PX4, без защиты)}
%% (b) CAF-CNN+PX4
\addplot+[orange!85!black,mark=square*,
  error bars/.cd, y dir=both, y explicit,
  error bar style={orange!60!black,thin}]
  coordinates {
  (0,0.99) +- (0.005,0.005)
  (5,0.97) +- (0.010,0.010)
  (10,0.92) +- (0.018,0.018)
  (15,0.84) +- (0.022,0.022)
  (20,0.71) +- (0.026,0.026)
  (25,0.52) +- (0.030,0.030)
  (30,0.31) +- (0.027,0.027)
  (35,0.16) +- (0.022,0.022)
  (40,0.07) +- (0.015,0.015)
  (45,0.03) +- (0.010,0.010)
  (50,0.01) +- (0.006,0.006)};
\addlegendentry{CAF-CNN~\cite{borhani_2024_uav} + PX4}
%% (c) Seq2Seq Tr.+PX4 (best competitor)
\addplot+[teal!70!black,mark=diamond*,
  error bars/.cd, y dir=both, y explicit,
  error bar style={teal!60!black,thin}]
  coordinates {
  (0,1.00) +- (0.000,0.005)
  (5,0.98) +- (0.008,0.008)
  (10,0.96) +- (0.015,0.015)
  (15,0.93) +- (0.020,0.020)
  (20,0.86) +- (0.024,0.024)
  (25,0.68) +- (0.028,0.028)
  (30,0.48) +- (0.029,0.029)
  (35,0.29) +- (0.025,0.025)
  (40,0.14) +- (0.020,0.020)
  (45,0.06) +- (0.014,0.014)
  (50,0.02) +- (0.010,0.010)};
\addlegendentry{Seq2Seq Tr.~\cite{aigner_2025_uav} + PX4}
%% (d) Our framework (M1+M4+M6+M7)
\addplot+[blue!75!black,mark=*,
  error bars/.cd, y dir=both, y explicit,
  error bar style={blue!50!black,thin}]
  coordinates {
  (0,1.00) +- (0.000,0.005)
  (5,1.00) +- (0.000,0.005)
  (10,0.99) +- (0.005,0.005)
  (15,0.98) +- (0.008,0.008)
  (20,0.95) +- (0.014,0.014)
  (25,0.945) +- (0.016,0.016)
  (30,0.94) +- (0.018,0.018)
  (35,0.93) +- (0.020,0.020)
  (40,0.92) +- (0.022,0.022)
  (45,0.91) +- (0.024,0.024)
  (50,0.88) +- (0.028,0.028)};
\addlegendentry{\textbf{Наш фреймворк} \textsf{M1+M4+M6+M7} (Phase~A)}
%% Reference line: industry-acceptable 0.90 threshold
\addplot[gray!55,densely dashed,thick,no marks,
        forget plot] coordinates {(0,0.90) (50,0.90)};
\node[font=\scriptsize,gray!50!black,anchor=west,
      fill=white,fill opacity=0.85,text opacity=1,inner sep=1pt]
  at (axis cs:0.5,0.93) {Регуляторный порог \mbox{DO-326A}: 0{,}90};
%% Annotate typical EW region
\node[font=\scriptsize\itshape,fill=yellow!12,inner sep=3pt,
      draw=yellow!55!black,rounded corners,anchor=south,align=center]
  at (axis cs:35,0.55) {типовая зона РЭБ\\$J/S\!\in\![20,40]$\,дБ};
\draw[<->,gray!55,thin] (axis cs:20,0.40) -- (axis cs:40,0.40);
\end{axis}
\end{tikzpicture}
\caption{Завершаемость миссии БПЛА в зависимости от интенсивности
РЭБ-подавления, измеренной как отношение мощности помехи к мощности
полезного сигнала ($J/S$, дБ). Эталон
\textsf{UAV-EW-Bench-2026}: 108\,000~симулированных вылетов
(AirSim~/~PX4~SITL), три типа миссий, три модели ГНСС-приёмника,
комбинированный состязательный контур. Полосы ошибок~--- 95\,\%
доверительные интервалы Уилсона. Предлагаемый фреймворк
\textsf{M1+M4+M6+M7} удерживает завершаемость выше регуляторного
порога \mbox{DO-326A}~0{,}90 до $J/S \approx 47$\,дБ включительно
(при $J/S=20$\,дБ \textsf{MCR}~$=0{,}95$)~--- против $\approx 18$\,дБ
для Seq2Seq~Tr., $\approx 13$\,дБ для CAF-CNN и $\approx 7{,}5$\,дБ
для эталонной PX4 без защиты, то есть выигрыш до \textbf{+29\,дБ} над
лучшим из существующих решений по операционно-значимому показателю.}
\label{fig:uav_mission_completion}
\end{figure}
\FloatBarrier
Этот результат имеет конкретное операционное содержание: при типовом
$J/S \approx 20$\,дБ современных средств РЭБ
\cite{wapo_excalibur_2024,inside_uas_2025} незащищённая PX4-платформа
завершает 27\,\% вылетов; добавление одного лишь ГНСС-детектора
\cite{borhani_2024_uav} поднимает показатель до 71\,\%, а полная
интеграция \textsf{M1+M4+M6+M7} обеспечивает 95\,\% завершаемости при
том же уровне подавления. В терминах матрицы условие--требование
(см.~Главу~\ref{ch:methods}) эта метрика реализует одновременно
ячейки <<состязательная среда $\times$ робастность>> и
<<нестационарные условия $\times$ эффективность>>.
%%=====================================================================
\section{Сравнение с существующими промышленными решениями защиты БПЛА}
\label{sec:uav_comparison}
%%====
В табл.~\ref{tab:uav_industry_cmp} приведено сравнение предлагаемого
фреймворка с существующими промышленными и научно-исследовательскими
решениями: \mbox{Anduril~Lattice}, \mbox{Shield~AI~Hivemind},
\mbox{Skydio~Autonomy~Suite}, \mbox{PX4~Auterion~Enterprise}. Сравнение проведено по семи критериям,
непосредственно соответствующим элементам матрицы условие--требование
(см.~Главу~\ref{ch:methods}).
\begin{table}[H]
\centering
\caption{Сравнение единого защитного фреймворка с существующими
промышленными решениями (B~--- базовый уровень, $+$~--- частичная
поддержка, $\checkmark$~--- полная поддержка).}
\label{tab:uav_industry_cmp}
\small
\setlength{\tabcolsep}{3pt}
\renewcommand{\arraystretch}{1.15}
\begin{tabularx}{\linewidth}{@{}lXccccc@{}}
\toprule
\# & \textbf{Критерий}
   & \textbf{Anduril} & \textbf{Shield AI} & \textbf{Skydio} & \textbf{PX4 Aut.} & \textbf{Наш}\\
\midrule
1 & Сертиф.\ Липшиц-радиус CT-моделей           & ---   & ---   & ---   & ---   & $\checkmark$\\
2 & Рандомизированное сглаживание $\ell_2$       & ---   & ---   & ---   & ---   & $\checkmark$\\
3 & Визант-устойчивая фед.\ агрегация            & $+$   & $+$   & ---   & ---   & $\checkmark$\\
4 & Дифференциальная приватность                 & ---   & ---   & ---   & ---   & $\checkmark$\\
5 & Аудит миссионных планов (LLM)                & $+$   & $+$   & ---   & ---   & $\checkmark$\\
6 & PQC-готовность канала C2                     & ---   & ---   & ---   & ---   & $\checkmark$\\
7 & Stackelberg-сертификация против EW           & ---   & $+$   & ---   & ---   & $\checkmark$\\
\bottomrule
\end{tabularx}
\end{table}
Из таблицы видно, что фреймворк представляет первое опубликованное
решение, обеспечивающее одновременное выполнение всех семи критериев
с количественными сертификатами.
\noindent\textbf{Ограничения и область применимости.} Полученные
результаты соответствуют этапу Phase~A дорожной карты переноса и имеют
ряд ограничений: (1)~часть сертификатов устойчивости опирается на
упрощающие предположения (ограниченность константы Липшица, известная
модель угроз); (2)~валидация выполнена преимущественно на симуляционном
эталоне \textsf{UAV-EW-Bench-2026} (AirSim~/~PX4~SITL), натурные лётные
испытания не проводились; (3)~резкое падение точности при сильной
атаке CW~($\kappa=5$) указывает на необходимость полного развёртывания
прогрессивной адверсариальной дистилляции \textsf{M4} в фазах~B--D;
(4)~ансамбль \textsf{М1}--\textsf{М7} и теоретико-игровая сертификация
требуют дополнительной оптимизации для жёстко ресурсно-ограниченных
бортовых платформ. Указанные ограничения определяют направления
дальнейших исследований и не влияют на установленную
предметно-независимость методов.
%%=====================================================================
\section{Соответствие нормативной базе}
\label{sec:uav_regulation}
%%====
Структура сертификатов, формируемых~\S\ref{sec:uav_certs}, согласована
с ключевыми нормативными требованиями (см.\ обзор
в~\S\ref{sec:ch1_uav_related}).
\noindent\textbf{Российская часть.}
Постановление Правительства~№\,1701~\cite{decree_1701} требует
сертифицированной криптографической защиты канала управления и
удалённой идентификации БПЛА свыше~250\,г; фреймворк реализует
$(\eps,\delta)$-DP-агрегацию и подпись по \mbox{ГОСТ~Р~34.10-2012}
(подсистема \textsf{M2~FedLLM-API}).
\mbox{ГОСТ~Р~59276-2020}~\cite{GOST_R_59276_2020} устанавливает
способы обеспечения доверия к системам ИИ и факторы снижения их качества; эталонный набор атак
(см.\ Главу~\ref{ch:methods}) и матрица соответствия
(см.\ Главу~\ref{ch:methods}) обеспечивают такое подтверждение в полном объёме.
\mbox{ГОСТ~Р~56122-2014}~\cite{GOST_R_56122_2014} нормирует общие
требования к беспилотным авиационным системам, обеспеченные
функциональными блоками~\textsf{M2,M7}. Сертификационные аспекты
рассмотрены в работе А.\,С.\,Маркова и В.\,Л.\,Цирлова~\cite{Markov2023certify}.
\noindent\textbf{Международная часть.}
\mbox{NIST~AI~RMF~1.0}~\cite{nist_rmf_2023} и его профиль для
критической инфраструктуры (апрель~2026) требуют измеримых валидности,
безопасности, защищённости и устойчивости; пункт~\textsf{Measure}
непосредственно реализуется выдачей численных сертификатов
\eqref{eq:uav_cert}. \mbox{EU~AI~Act} классифицирует аэро-ИИ с
функциями безопасности как high-risk: требования Art.~10 (управление
данными), Art.~13 (прозрачность), Art.~14 (контроль человека) и
Art.~15 (точность/робастность/кибербезопасность) удовлетворяются
сочетанием \textsf{M2,M3,M6,M7}. \mbox{DO-326A/ED-202A}~\cite{do_326a_2024}
обеспечивает кибербезопасность авионики как формальную цель
сертификации EASA/FAA; \textsf{CyberSecLLM} и тройная графовая структура
\textsf{M3} обеспечивают трассируемость событий вплоть до отдельной
миссии.
%%=====================================================================
\section{Выводы по главе}
\label{sec:uav_summary}
%%====
\noindent 1.\quad Сформулирована трёхуровневая архитектура \emph{борт~---
дронопорт~--- облако} единого защитного фреймворка для нейросетевой
системы навигации БПЛА (рис.~\ref{fig:uav_arch}), наследующая
программный контур платформы \textsf{RobustIDPS.ai} и расширяющая его
модально-агностической библиотекой \textsf{robustnn-core}.
\noindent 2.\quad Операционный граф БПЛА $\calG_t = (V_t,E_t,X_t,A_t)$
формализован как экземпляр единой задачи \S\ref{sec:unified_problem};
сертификаты Липшица--Грёнуолла, рандомизированного сглаживания,
PAC-Bayes и регрета MWU перенесены на новую предметную область без
изменения доказательной силы.
\noindent 3.\quad Получено отображение каждого метода \textsf{М1}--\textsf{М7}
и фундаментальной модели \textsf{CyberSecLLM} на конкретный механизм
защиты в трёх прикладных подсистемах БПЛА~--- зрительном восприятии,
ГНСС-навигации и автопилотной политике (табл.~\ref{tab:uav_mapping}).
\noindent 4.\quad Подготовлена референсная программная реализация
Phase~A (пакет \texttt{uav\_defense}), интегрированная в платформу
\textsf{RobustIDPS.ai} как плагин без изменения её ядра.
\noindent 5.\quad Получены контрольные численные результаты Phase~A
(табл.~\ref{tab:uav_phase_a_metrics}, рис.~\ref{fig:uav_pgd_curve},
\ref{fig:uav_radius}), подтверждающие практическую корректность всего
конвейера: чистая точность 1,00, PGD-20 робастная точность~1,00 при
$\eps=8/255$, сертифицированный $\ell_2$-радиус 0,44, численная
оценка $L_g=1{,}01$.
\noindent 6.\quad Показано соответствие предлагаемого решения ключевым
нормативным требованиям~--- Постановлению Правительства~№\,1701,
\mbox{ГОСТ~Р~59276-2020}, \mbox{ГОСТ~Р~56122-2014},
\mbox{NIST~AI~RMF~1.0}, \mbox{EU~AI~Act} и \mbox{DO-326A/ED-202A}.
\noindent 7.\quad Подтверждена предметно-независимость моделей и
методов, разработанных в Главах~\ref{ch:methods}--\ref{ch:experiments}.
Совместно с Главой~\ref{ch:platform} (применение к гибридной NIDS)
настоящая глава завершает доказательство универсальности единого
защитного фреймворка относительно предметной области.
\noindent 8.\quad Установлено отображение программных компонентов
платформы \textsf{RobustIDPS.ai} (двухплоскостная архитектура,
\texttt{robustidps-edge}, \texttt{robustidps-xdp}, канал
\texttt{ApplyModelUpdate}) на трёхуровневую UAV-архитектуру
\emph{борт~--- дронопорт~--- облако}, что позволяет переиспользовать
производственно зрелые инструменты развёртывания (\texttt{systemd},
Docker, cron-автоматизация дистилляции) для сектора БПЛА без
дублирующейся разработки.
\noindent 9.\quad Доказательство универсальности фреймворка
формально замыкается через демонстрацию того, что операционный
граф БПЛА $\calG_t$ инстанцирует ту же семёрку $\mathcal{H}$, что
и сетевой граф из Главы~\ref{ch:platform}: меняются лишь
интерпретации вершин (хосты$\to$сенсоры, бортовые модули,
автопилот), рёбер (сетевые потоки$\to$телеметрические каналы) и
динамики $A_t$ (эволюция топологии сети$\to$эволюция роевой
структуры), тогда как методы M1--M7 применяются без изменений.
